\section{Introduction}
\label{sec:introduction}

Mean reversion is a central organizing principle in continuous-time finance.
The Gaussian specification of \citet{Vasicek1977} and the square-root diffusion
of \citet{CoxIngersollRoss1985} remain canonical because they combine an
economically interpretable restoring force with tractable transition laws and,
in the interest-rate setting, tractable pricing implications.  Subsequent
empirical work shows, however, that the conditional distribution of financial
state variables is often more nonlinear and state dependent than either model
allows.  \citet{ChanEtAl1992} compare alternative short-rate diffusions,
\citet{AitSahalia1996} documents nonlinear drift and volatility, and
\citet{BibbySorensen1995}, \citet{AitSahalia2002}, and \citet{Kessler1997}
develop methods for discretely observed diffusions.  Standard treatments of
the underlying stochastic calculus and term-structure models include
\citet{Shreve2004} and \citet{BrigoMercurio2006}; the CIR boundary and solution
structure are further discussed by \citet{PlatenHeath2006}.  These results motivate a model class that
retains a precise notion of reversion while relaxing the tight link between a
diffusion's marginal law and its serial dependence.

Copulas provide a natural separation.  \citet{DarsowNguyenOlsen1992} relate
copulas to Markov transition operators through the star product, while
\citet{ChenFan2006} and \citet{ChenWuYi2009} develop semiparametric estimation
for stationary copula-based Markov models.  The diffusion construction of
\citet{BibbonaSacerdoteTorre2016} further shows how monotone space--time
transformations preserve serial dependence while changing the marginals, and
\citet{FangJiangLiuYang2020} studies copula-based Markov processes directly.
This literature establishes the probabilistic and econometric foundations for
modeling invariant marginals and temporal dependence separately.  It does not,
by itself, guarantee the directional restoring force required by an economic
notion of mean reversion.  Our theory makes that restriction explicit and
identifies when it is preserved by transformations and copula mixtures.

Stablecoins provide a particularly sharp application; \citet{AnteEtAl2023}
survey the rapidly growing empirical literature.  Their market value is
designed to remain close to a reference asset, so the relevant state variable
is the deviation from the peg rather than an unrestricted cryptocurrency
return.  Nevertheless, deviations are neither Gaussian nor uniformly short
lived.  \citet{LyonsViswanathNatraj2023} connect Tether's peg stability to the
design and accessibility of arbitrage, whereas \citet{DuanUrquhart2023} find
that the speed at which deviations are corrected varies substantially across
stablecoins and over time.  Time variation and nonlinear stabilization also
appear in \citet{DjogbenouInanJasiak2023} and
\citet{HuiWongLo2025}; \citet{BergaultEtAl2024} instead represent pegged-asset
dynamics with nested OU factors.  \citet{MaZengZhang2023} emphasizes the link
between run incentives and the centralization of arbitrage.  Tail events raise a separate run-risk dimension:
\citet{LiuMakarovSchoar2023} study the Terra--Luna collapse,
\citet{AhmedAldasoroDuley2025} show how information about reserve quality can
alter run incentives, and \citet{EichengreenNguyenViswanathNatraj2025}
construct market-based measures of stablecoin devaluation risk.  A credible
empirical model must therefore represent both ordinary correction toward the
peg and occasional transitions that are too wide, persistent, or asymmetric
for a single Gaussian diffusion.

This paper develops such a model in three steps.  Section~2 gives a
measure-theoretically precise Markov--copula foundation and distinguishes local
drift reversion from finite-horizon conditional-mean reversion and stationarity.
Section~3 constructs mean-reverting processes by monotone transformations of OU
and CIR diffusions.  The observable transition density follows from the inverse
map and its Jacobian, which also reveals when a formal differential-equation
solution fails as an empirical model because it is not globally monotone or
does not cover the observed support.  Section~4 introduces the central
copula-based construction.  Mixing a stationary transition copula with the
independence copula under a multiplicative survival weight yields a genuine
Markov semigroup, an exact transition density, and a Poisson-reset
interpretation.  The reset intensity augments the base mean-reversion rate
without altering the invariant marginal.

The empirical analysis evaluates these restrictions without selecting models
by complexity.  All theoretical families with a finite-dimensional observable
density are estimated or assigned a retained monotonicity, support, endpoint,
or identification diagnosis.  This is particularly important at five-minute
frequency, where measurement effects can distort OU estimates
\citep{HolyTomanova2018}.  Models are compared on a common log-price
conditional-density measure using an untouched 2024 validation sample, with
2025 retained as post-selection evidence.  This design matters for the main
conclusion.  The mixed-OU kernel produces a large in-sample likelihood gain and
the independent-reset component improves the validation score of Gaussian and
two-scale Gaussian copulas.  Yet it is redundant for the fitted Student-$t$
copula, deteriorates the fitted Archimedean mixtures, and does not by itself
reproduce the slow dependence observed across horizons.  A two-state OU obtains
the best five-minute predictive score, although its slow regime is weakly
identified.  Thus the copula-reset model is the paper's principal structural
contribution, while predictive model selection remains an empirical result
rather than a maintained claim.

The contribution is consequently both constructive and diagnostic.  On the
constructive side, the paper supplies verifiable conditions under which a
copula-based process has an observable Markov density and a restoring drift,
together with exact OU, CIR, reset, and pricing formulas.  On the diagnostic
side, it shows which formally derived transformations do not define usable
one-dimensional models and which empirical extensions require an augmented
state.  This distinction aligns the theoretical claims with the evidence and
prevents a better in-sample fit from being interpreted as universal
out-of-sample or economic identification.
